\begin{lemma}[Initial bridge]
Let $F$ be a front on $\mathbb N$ with $\emptyset\notin F$.  If
$m<n$, then
\[
  \mathsf{FrontBridge}\bigl(F,\{m\},\{n\},\{m,n\}\bigr).
\]
\end{lemma}
\begin{proof}
By \noderef{FrontTreeSingleton}, both singletons belong to $T(F)$.
The set $\{m,n\}$ is nonempty, and its least element is $m$, so its
finite bridge tail is $\{n\}$.  If $\{m\}\notin F$, then
$\{m\}$ is the proper initial segment of $\{m,n\}$ and the latter belongs
to $T(F)$: it is the ordered one-point extension supplied by
\noderef{FrontTreeOrderedExtension}.  If $\{m\}\in F$, the required left
initial-segment clause holds by equality.  The same two alternatives for
$\{n\}$ give either equality with the bridge tail or the terminal
initial-segment clause.  These are precisely all clauses of
\noderef{FrontBridge}, so the displayed bridge state holds.
\end{proof}
